\documentclass{report}

%%%%%%%%%%%%%%%%%%%%%%%%%%%%%%%%%
% PACKAGE IMPORTS
%%%%%%%%%%%%%%%%%%%%%%%%%%%%%%%%%


\usepackage[tmargin=2cm,rmargin=1in,lmargin=1in,margin=0.85in,bmargin=2cm,footskip=.2in]{geometry}
\usepackage{amsmath,lmodern,amsthm,amssymb,mathtools}
\usepackage[varbb]{newpxmath}
\usepackage{xfrac}
\usepackage[makeroom]{cancel}
\usepackage{mathtools}
\usepackage{bookmark}
\usepackage{enumitem}
\usepackage{hyperref,theoremref}
\hypersetup{
	pdftitle={Assignment},
	colorlinks=true, linkcolor=doc!90,
	bookmarksnumbered=true,
	bookmarksopen=true
}
\usepackage[most,many,breakable]{tcolorbox}
\usepackage{xcolor}
\usepackage{varwidth}
\usepackage{varwidth}
\usepackage{etoolbox}
%\usepackage{authblk}
\usepackage{nameref}
\usepackage{multicol,array}
\usepackage{tikz-cd}
\usepackage[ruled,vlined,linesnumbered]{algorithm2e}
\usepackage{comment} % enables the use of multi-line comments (\ifx \fi) 
\usepackage{import}
\usepackage{xifthen}
\usepackage{pdfpages}
\usepackage{transparent}

\newcommand\mycommfont[1]{\footnotesize\ttfamily\textcolor{blue}{#1}}
\SetCommentSty{mycommfont}
\newcommand{\incfig}[1]{%
    \def\svgwidth{\columnwidth}
    \import{./figures/}{#1.pdf_tex}
}

\usepackage{tikzsymbols}
\renewcommand\qedsymbol{$\Laughey$}



%%%%%%%%%%%%%%%%%%%%%%%%%%%%%%
% SELF MADE COLORS
%%%%%%%%%%%%%%%%%%%%%%%%%%%%%%



\definecolor{myg}{RGB}{56, 140, 70}
\definecolor{myb}{RGB}{45, 111, 177}
\definecolor{myr}{RGB}{199, 68, 64}
\definecolor{mytheorembg}{HTML}{F2F2F9}
\definecolor{mytheoremfr}{HTML}{00007B}
\definecolor{mylenmabg}{HTML}{FFFAF8}
\definecolor{mylenmafr}{HTML}{983b0f}
\definecolor{mypropbg}{HTML}{f2fbfc}
\definecolor{mypropfr}{HTML}{191971}
\definecolor{myexamplebg}{HTML}{F2FBF8}
\definecolor{myexamplefr}{HTML}{88D6D1}
\definecolor{myexampleti}{HTML}{2A7F7F}
\definecolor{mydefinitbg}{HTML}{E5E5FF}
\definecolor{mydefinitfr}{HTML}{3F3FA3}
\definecolor{notesgreen}{RGB}{0,162,0}
\definecolor{myp}{RGB}{197, 92, 212}
\definecolor{mygr}{HTML}{2C3338}
\definecolor{myred}{RGB}{127,0,0}
\definecolor{myyellow}{RGB}{169,121,69}
\definecolor{myexercisebg}{HTML}{F2FBF8}
\definecolor{myexercisefg}{HTML}{88D6D1}


%%%%%%%%%%%%%%%%%%%%%%%%%%%%
% TCOLORBOX SETUPS
%%%%%%%%%%%%%%%%%%%%%%%%%%%%

\setlength{\parindent}{1cm}


%================================
% DEFINITION BOX
%================================

\newtcbtheorem[number within=section]{Definition}{Definition}{enhanced,
	before skip=2mm,after skip=2mm, colback=red!5,colframe=red!80!black,boxrule=0.5mm,
	attach boxed title to top left={xshift=1cm,yshift*=1mm-\tcboxedtitleheight}, varwidth boxed title*=-3cm,
	boxed title style={frame code={
					\path[fill=tcbcolback]
					([yshift=-1mm,xshift=-1mm]frame.north west)
					arc[start angle=0,end angle=180,radius=1mm]
					([yshift=-1mm,xshift=1mm]frame.north east)
					arc[start angle=180,end angle=0,radius=1mm];
					\path[left color=tcbcolback!60!black,right color=tcbcolback!60!black,
						middle color=tcbcolback!80!black]
					([xshift=-2mm]frame.north west) \-- ([xshift=2mm]frame.north east)
					[rounded corners=1mm]-- ([xshift=1mm,yshift=-1mm]frame.north east)
					\-- (frame.south east) \-- (frame.south west)
					\-- ([xshift=-1mm,yshift=-1mm]frame.north west)
					[sharp corners]-- cycle;
				},interior engine=empty,
		},
	fonttitle=\bfseries,
	title={#2},#1}{def}

%================================
% ANSWER BOX
%================================

\newtcbtheorem[number within=section]{Answer}{Answer}{enhanced,
	before skip=2mm,after skip=2mm, colback=red!5,colframe=green!80!black,boxrule=0.5mm,
	attach boxed title to top left={xshift=1cm,yshift*=1mm-\tcboxedtitleheight}, varwidth boxed title*=-3cm,
	boxed title style={frame code={
					\path[fill=tcbcolback]
					([yshift=-1mm,xshift=-1mm]frame.north west)
					arc[start angle=0,end angle=180,radius=1mm]
					([yshift=-1mm,xshift=1mm]frame.north east)
					arc[start angle=180,end angle=0,radius=1mm];
					\path[left color=tcbcolback!60!black,right color=tcbcolback!60!black,
						middle color=tcbcolback!80!black]
					([xshift=-2mm]frame.north west) -- ([xshift=2mm]frame.north east)
					[rounded corners=1mm]-- ([xshift=1mm,yshift=-1mm]frame.north east)
					-- (frame.south east) -- (frame.south west)
					-- ([xshift=-1mm,yshift=-1mm]frame.north west)
					[sharp corners]-- cycle;
				},interior engine=empty,
		},
	fonttitle=\bfseries,
	title={#2},#1}{def}


%================================
% Solution BOX
%================================

\makeatletter
\newtcbtheorem{question}{Question}{enhanced,
	breakable,
	colback=white,
	colframe=myb!80!black,
	attach boxed title to top left={yshift*=-\tcboxedtitleheight},
	fonttitle=\bfseries,
	title={#2},
	boxed title size=title,
	boxed title style={%
			sharp corners,
			rounded corners=northwest,
			colback=tcbcolframe,
			boxrule=0pt,
		},
	underlay boxed title={%
			\path[fill=tcbcolframe] (title.south west)--(title.south east)
			to[out=0, in=180] ([xshift=5mm]title.east)--
			(title.center-|frame.east)
			[rounded corners=\kvtcb@arc] |-
			(frame.north) -| cycle;
		},
	#1
}{def}
\makeatother

%================================
% SOLUTION BOX
%================================

\makeatletter
\newtcolorbox{solution}{enhanced,
	breakable,
	colback=white,
	colframe=myg!80!black,
	attach boxed title to top left={yshift*=-\tcboxedtitleheight},
	title=Solution,
	boxed title size=title,
	boxed title style={%
			sharp corners,
			rounded corners=northwest,
			colback=tcbcolframe,
			boxrule=0pt,
		},
	underlay boxed title={%
			\path[fill=tcbcolframe] (title.south west)--(title.south east)
			to[out=0, in=180] ([xshift=5mm]title.east)--
			(title.center-|frame.east)
			[rounded corners=\kvtcb@arc] |-
			(frame.north) -| cycle;
		},
}
\makeatother

%================================
% Question BOX
%================================

\makeatletter
\newtcbtheorem{qstion}{Question}{enhanced,
	breakable,
	colback=white,
	colframe=mygr,
	attach boxed title to top left={yshift*=-\tcboxedtitleheight},
	fonttitle=\bfseries,
	title={#2},
	boxed title size=title,
	boxed title style={%
			sharp corners,
			rounded corners=northwest,
			colback=tcbcolframe,
			boxrule=0pt,
		},
	underlay boxed title={%
			\path[fill=tcbcolframe] (title.south west)--(title.south east)
			to[out=0, in=180] ([xshift=5mm]title.east)--
			(title.center-|frame.east)
			[rounded corners=\kvtcb@arc] |-
			(frame.north) -| cycle;
		},
	#1
}{def}
\makeatother

\newtcbtheorem[number within=chapter]{wconc}{Wrong Concept}{
	breakable,
	enhanced,
	colback=white,
	colframe=myr,
	arc=0pt,
	outer arc=0pt,
	fonttitle=\bfseries\sffamily\large,
	colbacktitle=myr,
	attach boxed title to top left={},
	boxed title style={
			enhanced,
			skin=enhancedfirst jigsaw,
			arc=3pt,
			bottom=0pt,
			interior style={fill=myr}
		},
	#1
}{def}


%%%%%%%%%%%%%%%%%%%%%%%%%%%%%%
% SELF MADE COMMANDS
%%%%%%%%%%%%%%%%%%%%%%%%%%%%%%

\newcommand{\ans}[2]{\begin{Answer}[colbacktitle=green!75!black]{#1}{}#2\end{Answer}}
\newcommand{\dfn}[2]{\begin{Definition}[colbacktitle=red!75!black]{#1}{}#2\end{Definition}}
\newcommand{\dfnc}[2]{\begin{definition}[colbacktitle=red!75!black]{#1}{}#2\end{definition}}
\newcommand{\qs}[2]{\begin{question}{#1}{}#2\end{question}}


%%%%%%%%%%%%%%%%%%%%%%%%%%%%%%%%%%%%%%%%%%%
% TABLE OF CONTENTS
%%%%%%%%%%%%%%%%%%%%%%%%%%%%%%%%%%%%%%%%%%%

\usepackage{tikz}
\definecolor{doc}{RGB}{0,60,110}
\usepackage{titletoc}
\contentsmargin{0cm}
\titlecontents{chapter}[3.7pc]
{\addvspace{30pt}%
	\begin{tikzpicture}[remember picture, overlay]%
		\draw[fill=doc!60,draw=doc!60] (-7,-.1) rectangle (-0.9,.5);%
		\pgftext[left,x=-3.5cm,y=0.2cm]{\color{white}\Large\sc\bfseries Chapter\ \thecontentslabel};%
	\end{tikzpicture}\color{doc!60}\large\sc\bfseries}%
{}
{}
{\;\titlerule\;\large\sc\bfseries Page \thecontentspage
	\begin{tikzpicture}[remember picture, overlay]
		\draw[fill=doc!60,draw=doc!60] (2pt,0) rectangle (4,0.1pt);
	\end{tikzpicture}}%
\titlecontents{section}[3.7pc]
{\addvspace{2pt}}
{\contentslabel[\thecontentslabel]{2pc}}
{}
{\hfill\small \thecontentspage}
[]
\titlecontents*{subsection}[3.7pc]
{\addvspace{-1pt}\small}
{}
{}
{\ --- \small\thecontentspage}
[ \textbullet\ ][]

\makeatletter
\renewcommand{\tableofcontents}{%
	\chapter*{%
	  \vspace*{-20\p@}%
	  \begin{tikzpicture}[remember picture, overlay]%
		  \pgftext[right,x=15cm,y=0.2cm]{\color{doc!60}\Huge\sc\bfseries \contentsname};%
		  \draw[fill=doc!60,draw=doc!60] (13,-.75) rectangle (20,1);%
		  \clip (13,-.75) rectangle (20,1);
		  \pgftext[right,x=15cm,y=0.2cm]{\color{white}\Huge\sc\bfseries \contentsname};%
	  \end{tikzpicture}}%
	\@starttoc{toc}}
\makeatother

% %From M275 "Topology" at SJSU
\newcommand{\id}{\mathrm{id}}
\newcommand{\taking}[1]{\xrightarrow{#1}}
\newcommand{\inv}{^{-1}}

%From M170 "Introduction to Graph Theory" at SJSU
\DeclareMathOperator{\diam}{diam}
\DeclareMathOperator{\ord}{ord}
\newcommand{\defeq}{\overset{\mathrm{def}}{=}}

%From the USAMO .tex files
\newcommand{\ts}{\textsuperscript}
\newcommand{\dg}{^\circ}
\newcommand{\ii}{\item}

% % From Math 55 and Math 145 at Harvard
% \newenvironment{subproof}[1][Proof]{%
% \begin{proof}[#1] \renewcommand{\qedsymbol}{$\blacksquare$}}%
% {\end{proof}}

\newcommand{\liff}{\leftrightarrow}
\newcommand{\lthen}{\rightarrow}
\newcommand{\opname}{\operatorname}
\newcommand{\surjto}{\twoheadrightarrow}
\newcommand{\injto}{\hookrightarrow}
\newcommand{\On}{\mathrm{On}} % ordinals
\DeclareMathOperator{\img}{im} % Image
\DeclareMathOperator{\Img}{Im} % Image
\DeclareMathOperator{\coker}{coker} % Cokernel
\DeclareMathOperator{\Coker}{Coker} % Cokernel
\DeclareMathOperator{\Ker}{Ker} % Kernel
\DeclareMathOperator{\rank}{rank}
\DeclareMathOperator{\Spec}{Spec} % spectrum
\DeclareMathOperator{\Tr}{Tr} % trace
\DeclareMathOperator{\pr}{pr} % projection
\DeclareMathOperator{\ext}{ext} % extension
\DeclareMathOperator{\pred}{pred} % predecessor
\DeclareMathOperator{\dom}{dom} % domain
\DeclareMathOperator{\ran}{ran} % range
\DeclareMathOperator{\Hom}{Hom} % homomorphism
\DeclareMathOperator{\Mor}{Mor} % morphisms
\DeclareMathOperator{\End}{End} % endomorphism



\newcommand{\hl}[2][grayCF]{%
    \tcbhighmath[colback=#1,%!30!white,
                boxrule=0pt,
                arc=2pt,
                left=1pt,right=1pt,top=.1pt,bottom=.1pt]{#2}
                % outer arc=2pt]{#2}
}
\newtcbox{\hlt}[1][gray!20]{%
  on line,
  colback=#1,
  colframe=#1,
  boxrule=0pt,
  arc=2pt,
  left=2pt,
  right=2pt,
  top=1pt,
  bottom=1pt,
  boxsep=0pt,
  verbatim
}




% ICONS
\newcommand{\icons}[1]{\includegraphics[height=1.2em]{icons/#1}}

% Specific aliases
\newcommand{\engine}{\icons{engine.png}}
\newcommand{\file}{\icons{file.png}}
\newcommand{\folder}{\icons{folder.png}}
\newcommand{\game}{\icons{game.png}}
\newcommand{\gear}{\icons{gear.png}}
\newcommand{\calc}{\icons{math.png}}
\newcommand{\package}{\icons{package.png}}
\newcommand{\greenCheck}{\icons{greenCheck.png}}

\begin{document}

\chapter{Mission Bus Integration Lab (MBIL)}
% ==============================================================================================================
% 1 
% ==============================================================================================================
\section{Project purpose}
\vspace{\baselineskip}
Mission Bus Integration Lab (MBIL) is a deterministic, configuration-driven software system that simulates a simplified aerospace mission environment built around a central data bus, redundant mission computers, subsystem messaging, defined failure modes, observability, and failover behavior.
\vspace{\baselineskip}

The goal is not to make a flashy simulator
\vspace{\baselineskip}

The goal to demonstrate that I understand how \textbf{mission-critical systems behave, fail, reover, and get debugged.}

\begin{center}
\rule{0.9\textwidth}{1pt}
\end{center}
% ==============================================================================================================
% 2 
% ==============================================================================================================

\section{What MBIL should communicate}
\vspace{\baselineskip}
When someone sees this project, they should think:
\begin{itemize}
   \item[$\star$] This person understands relibable system design. 
   \item This person thinks in interfaces and contracts
   \item This person knows how to model faults and degraded operation
   \item This person knows how to make systems observable
   \item This person can separate control logic, transport, config, and monitoring
\end{itemize}
MBIL should read like a \textbf{systems engineering lab}, not a game and not a toy simulation 

\begin{center}
\rule{0.9\textwidth}{1pt}
\end{center}
% ==============================================================================================================
% 3 
% ==============================================================================================================

\section{Core project statement}
MBIL is a \textbf{tick-driven mission-system simulation platform} with these defining properties:
\begin{itemize}
   \item Deterministic execution
   \item Configuration-defined topology
   \item Redundant control nodes
   \item First-class failure scenarios
   \item structured observability 
   \item Recoverable degraded operation 
\end{itemize}

\begin{center}
\rule{0.9\textwidth}{1pt}
\end{center}
% ==============================================================================================================
% 4 
% ==============================================================================================================
\section{Design priorities}
These are the real priorities, in order:
\begin{itemize}
   \item[] \textbf{Priority 1 -- Determinism}
\vspace{\baselineskip}The same scenario should produce the same behavior every run unless randomness is explicitly enabled 
   \item[] \textbf{Priority 2 -- Reliability modeling}
\vspace{\baselineskip}Failures are not side noes.  They are part of the system design. 
   \item[] \textbf{Priority 3 -- Observability}
\vspace{\baselineskip}Every important decision, state transition, and message path should be visible in logs or snapshots. 
   \item[] \textbf{Priority 4 -- Clear modular architecture}
\vspace{\baselineskip}Subsystem logic, bus behavior, timing, fault injection, and loggin must stay cleanly separated. 
   \item[] \textbf{Priority 5 -- Configurability}
\vspace{\baselineskip}System topology and scenarios should be changel from config files, not by rewriting core logic. 
\end{itemize}

\begin{center}
\rule{0.9\textwidth}{1pt}
\end{center}
% ==============================================================================================================
% 5 
% ==============================================================================================================
\section{Function Goals}
MBIL must support:
\begin{itemize}
   \item Subsystem-to-subsystem communication through a central bus 
   \item Multiple sybsystem types 
   \item Dual mission computer architecture 
   \item Heartbeat monitoring 
   \item Failover behavior 
   \item Defined failure injection 
   \item Message routing
   \item Deterministic Scheduling 
   \item Structured logs 
   \item State Snapshots 
   \item Scenario-based runs
\end{itemize}

\begin{center}
\rule{0.9\textwidth}{1pt}
\end{center}
% ==============================================================================================================
% 6 
% ==============================================================================================================
\section{High-level Architecture}
MBIL should be split into these major layers. 
\vspace{\baselineskip} \textbf{A. Simulation Core}

\textbf{Responsible for:}
\begin{itemize}
   \item Tick loop 
   \item Execution order 
   \item Simulation Duration 
   \item System lifecycle 
   \item Calling each subsystem in deterministic order
\end{itemize}
\textbf{B. Bus Layer}
\textbf{Responsible for:}
\begin{itemize}
   \item Receiving Outgoing messages 
   \item Storing pending messages 
   \item Applying routing rules 
   \item Applying message-related faults 
   \item Delivering ready messages 
   \item Preserving deterministic ordering
\end{itemize}
\vspace{\baselineskip} \textbf{C. Subsystem Layer}
\textbf{Responsible for:}
\begin{itemize}
   \item Mission Computers 
   \item Sensors 
   \item Display / Endpoints
   \item Health State 
   \item Behavior per tick
   \item Message Handling
\end{itemize}
\vspace{\baselineskip} \textbf{D. Fault Management Layer}
\textbf{Responsible for:}
\begin{itemize}
   \item Defined Failure modes
   \item Scripted Scenario Events
   \item Optional Random Fault behavior later
   \item Recording Fault Activations
   \item Applying Degradation Cleanly
\end{itemize}
\vspace{\baselineskip} \textbf{E. Observability Layer}
\textbf{Responsible for:}
\begin{itemize}
   \item Structured Event Logging
   \item Traceable Messgage Lifecycle
   \item State Snapshots
   \item Failover Reportin
   \item Auditability
\end{itemize}
\vspace{\baselineskip} \textbf{F. Configuration Layer}
\textbf{Responsible for:}
\begin{itemize}
   \item Loading JSON config
   \item Validating Topology
   \item Instantiating Subsystem Objects
   \item Loading Scenario/Fault Definitions
   \item Setting timing and thresholds
\end{itemize}

\begin{center}
\rule{0.9\textwidth}{1pt}
\end{center}
% ==============================================================================================================
% 7 
% ==============================================================================================================
\section{Runtime Model}
MBIL should use a \textbf{fixed tick} simulation model. 

\textbf{Example:}
\begin{itemize}
   \item One tick every 100 ms simulated time 
   \item Run for 300 ticks per scenario
\end{itemize}

\textbf{This is important because it gives you:}
\begin{itemize}
   \item Deterministic Behavior
   \item Clean Testability
   \item Easier Fault scripting
   \item Easier replay and debugging
\end{itemize}

\textbf{Tick Processing Order}

Use a Fixed order every tick:

\textbf{Phase 1 -- Scenario updates}
\begin{itemize}
   \item Activate any scheduled Failures
   \item Deactivate expired Failrues 
   \item Apply Manual state changes if configured 
\end{itemize}

\textbf{Phase 2 -- Sensor Updates}
\begin{itemize}
   \item Sensors decide whether to emit new data this tick
\end{itemize}

\textbf{Phase 3 -- Mission Computer Updates}
\begin{itemize}
   \item Active controller processes a state 
   \item Standby Controller monitors health 
   \item Heartbeats generated if healthy
\end{itemize}

\textbf{Phase 4 -- Endpoint / Display Updates}
\begin{itemize}
   \item Display or logger consumers update their local state. 
\end{itemize}

\textbf{Phase 5 -- Bus Processing}
\begin{itemize}
   \item Pending messages sorted and processsed
   \item Faults applied to eligible messages
   \item Delayed messages retained
   \item Deliverable messages routed
\end{itemize}

\textbf{Phase 6 -- Observability Output}
\begin{itemize}
   \item Event logs Written
   \item Snapshot Optionally captured
\end{itemize}

That order should be documented and never casually changed.

\begin{center}
\rule{0.9\textwidth}{1pt}
\end{center}
% ==============================================================================================================
% 8 
% ==============================================================================================================
\section{Deterministic behavior rules}

This is now a first-class design feature. 

\vspace{1\baselineskip}
\textbf{MBIL determinism rules}
\vspace{2\baselineskip}
\vspace{\baselineskip}\textbf{Rule 1: fixed tick rate}
\vspace{\baselineskip}\vspace{1\baselineskip}The simulation advances in discrete ticks. 

\noindent \textbf{Rule 2: stable subsystem order}
\vspace{\baselineskip}\vspace{1\baselineskip}Subsystems are always ticked in the same deterministic order, such as sorted by ID. 

\noindent \textbf{Rule 3: Stable message ordering}
\vspace{\baselineskip}\vspace{1\baselineskip}Messages created in the same tick are delivered in a defined order:
\begin{itemize}
   \item Priority 
   \item Tick created 
   \item Source ID 
   \item Message sequence number
\end{itemize}

\noindent \textbf{Rule 4: No hidden asynchronous behavior in MVP}
\vspace{\baselineskip}\vspace{1\baselineskip}No real multithreading in the initial version. 

\noindent \textbf{Rule 5: Reproducible scenario results}
\vspace{\baselineskip}\vspace{1\baselineskip}A scenario file should produce identical output unless randomness is explicitly enabled with a fixed seed.  
\vspace{\baselineskip}\vspace{1\baselineskip}This matters because it makes the system feel like an engineering lab rather than an app demo. 

\begin{center}
\rule{0.9\textwidth}{1pt}
\end{center}
% ==============================================================================================================
% 9 
% ==============================================================================================================
\section{Core subsystem model}

Each subsystem should implement a common interface. 

\begin{itemize}
   \item id()
   \item type()
   \item tick(context)
   \item receive(message)
   \item health()
   \item snapshot()
\end{itemize}

\textbf{Health states}

Every subsystem should have explicit health states such as:
\begin{itemize}
   \item Nominal
   \item Degraded
   \item Offline
   \item Failed
   \item Recovering
\end{itemize}

This makes failure hanling more professional and easier to test.

\begin{center}
\rule{0.9\textwidth}{1pt}
\end{center}
% ==============================================================================================================
% 10 
% ==============================================================================================================
\section{Main subsystem types}

\textbf{A. Mission Computer}

\textbf{Two instances:}
\begin{itemize}
   \item \textbf{MC1} = primary
   \item \textbf{MC2} = backup
\end{itemize}

\textbf{Responsibilities:}
\begin{itemize}
   \item Process incoming sensor/status messages
   \item Emit heartbeats
   \item Generate command/status output
   \item Maintain own health state
   \item Participate in failover logic
\end{itemize}


\textbf{MC1 Behavior}
\begin{itemize}
   \item Starts as active controller
   \item Processes sensor inputs
   \item Emits heartbeat periodically
   \item Produces system status and commands
\end{itemize}

\textbf{MC2 Behavior}
\begin{itemize}
   \item Starts in standby
   \item Watches for MC1 heartbeat loss
   \item Tracks timeout threshold
   \item Assumes control when failover condition is met 
   \item Announces failover in log/event stream
\end{itemize}

\begin{center}
\rule{0.9\textwidth}{1pt}
\end{center}

\textbf{B. Sensors}

\textbf{Example Sensors} 
\begin{itemize}
   \item Altitude Sensor
   \item Temperature Sensor
   \item Navigation Source
\end{itemize}

\textbf{Responsibilities} 
\begin{itemize}
   \item Emit data on configured interval
   \item Send structured sensor messages
   \item Support degraded or failed behavior under fault scenarios
\end{itemize}

\textbf{Sensor outputs should be simple but believable:} 
\begin{itemize}
   \item Value
   \item Valid Flag
   \item Quality / State
\end{itemize}


\begin{center}
\rule{0.9\textwidth}{1pt}
\end{center}


\textbf{C. Display/Endpoint}

\textbf{Responsibilities:} 
\begin{itemize}
   \item Receive commands/status
   \item Display system state through logs or summary output
   \item Report currently active mission computer
   \item Show degraded conditions
\end{itemize}

This is not a GUI priority.  Console or structured log output is engouth.

\begin{center}
\rule{0.9\textwidth}{1pt}
\end{center}
% ==============================================================================================================
% 11 
% ==============================================================================================================
\section{Message bus design}

The bus is central, but should stay simple. 

\textbf{Bus responsibilities}
\begin{itemize}
   \item Accept outgoing messages from subsystems
   \item Assign message metadata if needed
   \item Hold pending and delayed queues
   \item Apply failure rules
   \item Route messages by destination
   \item Support broadcast
   \item log delivery outcome
\end{itemize}

\textbf{Message routing modes}

\vspace{\baselineskip} Support:
\begin{itemize}
   \item point-to-point
   \item broadcast
   \item optional group routing later
\end{itemize}


\begin{center}
\rule{0.9\textwidth}{1pt}
\end{center}
% ==============================================================================================================
% 12 
% ==============================================================================================================
\section{Message model}

\vspace{\baselineskip} Keep the message model explicit and interview-friendly. 

\noindent\textbf{Message fields}


\noindent Each message should include:
\begin{itemize}
   \item message\_id
   \item sequence\_number
   \item tick\_created
   \item tick\_due
   \item source\_id
   \item destination\_id
   \item message\_type
   \item priority
   \item payload
   \item trace\_id or correlation ID
\end{itemize}
\vspace{\baselineskip} That last field is very good for observability. 

\textbf{Message types}

Starts with:
\begin{itemize}
   \item HEARTBEAT
   \item SENSOR\_DATA
   \item STATUS
   \item COMMAND
   \item ALERT
   \item FAILOVER\_NOTICE
\end{itemize}

\textbf{Payload model}

Use a simple tagged payload design. 

Examples:
\begin{itemize}
   \item Heartbeat payload: current role, heatlh
   \item Sensor payload: sensor name, numeric value, validity
   \item Status payload: active controller, mode, health summary
   \item Command payload: command name, parameter
\end{itemize}

\begin{center}
\rule{0.9\textwidth}{1pt}
\end{center}
% ==============================================================================================================
% 13 
% ==============================================================================================================
\section{Failure modes as a first-class feature}

This is one of the the biggest improvements. 

Do not describ ehtis as "supports faults."

Describe it as:

\textbf{MBIL includes defined operational failure modes and corresponding system responses.}

\textbf{Core failure mdoe categories}

\textbf{1. Bus timeout / delivery failure}

A Message does not arrive in expected time.  

Effect:
\begin{itemize}
   \item Receiving system may mark source stale
   \item Failover logic may be triggered if heartbeats are affected
\end{itemize}

\textbf{3. Message delay}

Message is delivered after configured delay. 

Effect:
\begin{itemize}
   \item Stale data
   \item Delayed heartbeat detection
   \item Timing edge cases
\end{itemize}

\textbf{4. Message corruption}

Payload kared invalid or altered. 

Effect:
\begin{itemize}
   \item Receiving system rejects or flags data
   \item Degraded mode may be entered
\end{itemize}

\textbf{5. Node failure}

Subsystem goes offline. 

Examples:
\begin{itemize}
   \item MC1 failes
   \item Sensor fails
   \item Display goes offline
\end{itemize}

\textbf{6. Node degradation}

Subsystem still runs but performs poorly. 

Examples:
\begin{itemize}
   \item Lower sensor update rate
   \item Heartbeat jitter
   \item Intermittent bad data
\end{itemize}


\begin{center}
\rule{0.9\textwidth}{1pt}
\end{center}
% ==============================================================================================================
% 14 
% ==============================================================================================================



\end{document}
